% Appendix chapter — populated with thesis-relevant material.

\noindent The appendices collect four artefacts that support the main
text but would interrupt its narrative: the memory-bus HTTP contract
(Appendix~A), the canonical long-format results-dataframe schema used by
every experiment runner (Appendix~B), an example C1 workload
instance from family~(a) showing the source-fragment structure
(Appendix~C), and a one-command reproducibility recipe together with
an illustrative \texttt{curl} trace (with timestamps and chain hashes
elided for legibility) exercising the memory-bus endpoints
(Appendix~D). The full reproducibility package, including the
model cards and the data cards, is shipped with the reference
release referenced in the closing section of
\charef{summary}.

\section{Appendix A. Memory-Bus HTTP Contract}
\markboth{left}{\normalfont{Appendix A. Memory-Bus HTTP Contract}}

The reference memory-bus service exposes four endpoints. The contract
below summarises the request and response shapes; the authoritative
OpenAPI schema lives at \texttt{docs/openapi.json} in the repository
and is regenerated via \texttt{make bus-openapi}.

\begin{description}\setlength\itemsep{0pt}
\item[\texttt{POST /v1/write}] Body: a \texttt{Fragment} JSON object
(\texttt{fragment\_id}, \texttt{text}, \texttt{tags}) and an optional
\texttt{target\_ratio}. Returns the assigned \texttt{slot\_id}, the
audit row identifier, and the hex-encoded \texttt{chain\_hash}.
HTTP $201$ on success; $403$ on policy denial.

\item[\texttt{GET /v1/read/\{slot\_id\}}] Returns a
\texttt{CompressedSlot} with its payload, tags, compressor identifier,
nominal ratio, and creation timestamp. HTTP $200$ on success;
$403$ on policy denial; $404$ on slot not found.

\item[\texttt{POST /v1/subscribe}] Body: a \texttt{SubscribeRequest}
with a \texttt{query} string, a \texttt{ttl\_seconds} positive integer,
and a top-\textit{k} integer. Returns a server-sent-event stream
emitting newly matched \texttt{slot\_id}s with their cosine scores
until the TTL elapses.

\item[\texttt{GET /v1/audit/\{slot\_id\}}] Returns the chronological
list of \texttt{AuditRow} entries that mention the given slot.
\end{description}

Every request is processed by the \texttt{PolicyMiddleware}, which
parses the requester principal from the \texttt{X-M6-Principal} header
in the development build and from an OAuth/JWT verifier in a production
deployment. The principal carries a $64$-bit access-control mask and a
$5$-tier classification level (PUBLIC $<$ INTERNAL $<$ CONFIDENTIAL $<$
RESTRICTED $<$ SECRET). A request is granted access to a slot when the
principal's mask is a superset of the slot tag's mask and the
principal's classification is at least the slot tag's classification.

\section{Appendix B. Long-Format Results Schema}
\markboth{left}{\normalfont{Appendix B. Long-Format Results Schema}}

Every experiment runner writes its results to a CSV that conforms to
the long-format schema implemented as
\texttt{m6.evaluation.statistics.LongDF}. The columns are listed in
Table~\ref{tab:longdf}.

\begin{table}[h]
\centering
\caption{The canonical long-format dataframe schema used by every
experiment runner. The same schema underlies every CSV under
\texttt{results/}, which makes cross-hypothesis analysis a single
\texttt{pandas.concat} away.}
\label{tab:longdf}
\small
\begin{tabular}{p{4cm}p{2.5cm}p{6.5cm}}
\hline
Column & Type & Notes \\
\hline
\texttt{experiment\_id} & str & unique per run \\
\texttt{hypothesis} & str & one of \texttt{h1}\,\ldots\,\texttt{h6} or \texttt{caac}, \texttt{frontier}, \texttt{hotpotqa} \\
\texttt{compressor} & str & e.g.\ \texttt{lingua2}, \texttt{filter},
\texttt{phi3-extractive}, \texttt{truncation}, \texttt{caac}, \texttt{identity} \\
\texttt{ratio} & float & nominal compression ratio \\
\texttt{actual\_ratio} & float & measured ratio post-compression \\
\texttt{pipeline} & str & one of \texttt{none}, \texttt{P1},
\texttt{P2}, \texttt{P3} \\
\texttt{model} & str & e.g.\ \texttt{llama-3.1-8b-instruct} \\
\texttt{model\_size} & str & \texttt{1.5b}, \texttt{3.8b}, \texttt{8b}, \texttt{72b} \\
\texttt{workload\_family} & str & \texttt{a}, \texttt{b}, or \texttt{c} \\
\texttt{workload\_id} & str & e.g.\ \texttt{c1-a-007} \\
\texttt{seed} & int &  \\
\texttt{metric} & str & e.g.\ \texttt{coord\_success}, \texttt{qa\_f1},
\texttt{preservation\_rate} \\
\texttt{value} & float &  \\
\texttt{wallclock\_ms} & int &  \\
\texttt{eur\_cost} & float &  \\
\texttt{run\_id} & str &  \\
\texttt{git\_sha} & str &  \\
\texttt{config\_hash} & str &  \\
\texttt{created\_at} & str & ISO-8601 UTC \\
\texttt{invalid} & bool & control-condition variance dominates? \\
\texttt{invalid\_reason} & str & free-text explanation if invalid \\
\hline
\end{tabular}
\end{table}

\section{Appendix C. Example C1 Workload (family a)}
\markboth{left}{\normalfont{Appendix C. Example C1 Workload}}

The following JSON object illustrates a single instance of C1
family~(a) (cross-document fact aggregation). Each fragment represents
one institutional system; the planner is asked to aggregate hours and
budget across all eight; the critic verifies the answer against the
ground-truth aggregate.

\begin{verbatim}
{
  "workload_id": "c1-a-007",
  "family": "a",
  "seed": 1247,
  "tag_distribution": "skewed",
  "initial_prompt":
    "Aggregate the period-end report for project P-0001 across the
     following systems: system_1, system_2, system_3, system_4,
     system_5, system_6, system_7, system_8. Report total hours and total
     budget.",
  "expected_answer": "hours=853;budget=148340",
  "n_agents": 8,
  "fragments": [
    {
      "fragment_id": "c1-a-007/system_1",
      "text": "[system_1] Project P-0001 - period 2025-Q3. ...",
      "tags": {"acl_mask": 12876, "classification": 1}
    },
    {
      "fragment_id": "c1-a-007/system_2",
      "text": "[system_2] Project P-0001 - period 2025-Q3. ...",
      "tags": {"acl_mask": 8901, "classification": 2}
    }
    // (six more fragments, one per system)
  ],
  "sub_tasks": [
    {
      "sub_task_id": "c1-a-007/sub-0",
      "description":
        "Extract recorded hours and budget for project P-0001 from
         system_1.",
      "expected_solver": "worker-0",
      "expected_answer": "hours=112;budget=18420"
    }
    // (seven more sub-tasks)
  ],
  "protected_facts": [
    {
      "fragment_id": "c1-a-007/system_2",
      "fact": "system_2.budget=22150",
      "yesno_questions": [
        "Did system_2 for project P-0001 exceed EUR 27150 in budget?",
        "Was system_2's recorded hours for P-0001 more than 142?"
      ],
      "answers": ["no", "no"]
    }
  ]
}
\end{verbatim}

The full dataset is reproduced from
\texttt{configs/benchmark/c1-v0.1.yaml} via \texttt{make
bench-generate}; the manifest carries the configuration hash so any
reader can verify they are looking at the same release.

\section{Appendix D. Reproducibility Recipe and Memory-Bus Trace}
\markboth{left}{\normalfont{Appendix D. Reproducibility Recipe and Memory-Bus Trace}}

\subsection{D.1 One-command reproduction of every chapter figure}

The reproducibility package is the complete reference repository
released to examiners and collaborators.
Historical planning documents in the repository reflect earlier
protocol designs; the canonical experiment protocol for the submitted
thesis is the final set of experiment runners and result artefacts
listed here.
The package contains a \texttt{docker-compose.yml} that
brings the memory-bus reference service up locally, a GitHub
release tag matched to the manuscript, model and data cards under
\texttt{docs/}, and a \texttt{README.md} that exposes
one-command reproduction targets for every headline figure (the
compute envelope and software stack are detailed in
Section~D.3):

\begin{description}\setlength\itemsep{1pt}
\item[\texttt{make repro-bench}]
regenerates the $150$-instance C1 benchmark from
\texttt{configs/benchmark/c1-v0.1.yaml}. Output:
\texttt{data/processed/c1-v0.1/}.
\item[\texttt{make repro-cache}]
runs the compression precomputation for all four compressors at
the ten canonical ratios across the $150$ workloads. Requires
the GPU host; total wallclock $\sim 14$~hours. Output:
\texttt{results/compression\_cache/}.
\item[\texttt{make repro-cliff}]
runs the H1/H2 cliff sweep against the compression cache,
producing all of \charef{experiments}'s Section~\ref{sec:h1}
through Section~\ref{sec:cem_match} figures. Wallclock
$\sim 3$~hours on the M4~Pro.
\item[\texttt{make repro-scaling}]
runs the model-scaling sweep at $1.5$B, $3.8$B, and $8$B with the
local Ollama backend, producing Corollary~1's figures. Wallclock
$\sim 90$~minutes.
\item[\texttt{make repro-frontier}]
runs the cliff sweep against the Featherless API for
Qwen-2.5-72B-Instruct and DeepSeek~V4~Pro, the two
calibrated-regime frontier validators. (The out-of-regime
GPT-oss-120B diagnostic is scoped out of the canonical recipe as
an extended-reasoning planner and is not part of this target.) Requires a
\texttt{FEATHERLESS\_API\_KEY}. Output:
\texttt{results/frontier\_\{qwen72b,deepseekv4\}/}.
\item[\texttt{make repro-rag}]
runs the H3 RAG pipeline catalogue, producing
Section~\ref{sec:h3}'s figures. Wallclock $\sim 3$~hours.
\item[\texttt{make repro-disclosure}]
runs the H4 inference-disclosure measurement against the unbiased
benchmark, producing Section~\ref{sec:h4_disclosure}'s figures.
Wallclock $\sim 1$~hour on the M4~Pro.
\item[\texttt{make repro-transfer}]
runs the HotpotQA and MultiHop-RAG transfer arm, producing
Section~\ref{sec:corollary2}'s figures. Wallclock $\sim 2$~hours.
\item[\texttt{make repro-figs}]
regenerates every PDF and PNG figure under \texttt{figures/} from
the canonical CSVs. Pure-Python; runs in under a minute.
\end{description}

The \texttt{make repro-all} target chains the six laptop-runnable
steps above (\texttt{repro-cliff}, \texttt{repro-scaling},
\texttt{repro-rag}, \texttt{repro-disclosure},
\texttt{repro-transfer}, and \texttt{repro-figs}) in dependency
order, excluding the two that need special resources:
\texttt{repro-cache} (GPU host) and \texttt{repro-frontier} (a
Featherless API key), which must be run explicitly; the whole
package is $\sim 30$~hours of GPU time and $\sim 12$~hours of laptop
time end-to-end. The byte-deterministic / in-distribution /
not-reproducible status of each target follows the three-tier
guarantee of Section~\ref{sec:repro} (\charef{experiments}):
\texttt{repro-cache}, \texttt{-cliff}, \texttt{-rag},
\texttt{-disclosure}, and \texttt{-figs} are Tier~1,
\texttt{repro-scaling} is Tier~2, and \texttt{repro-frontier} is
Tier~3.

\subsection{D.2 Memory-bus end-to-end trace}

The following \texttt{curl} sequence exercises the four
endpoints of the memory bus end-to-end against the
\texttt{docker-compose}-up'ed reference service. The trace is
the rubric-\#5 evidence for the C4 contribution: the memory bus
is a running system, not a paper design. The output of each step
is reproduced from a run against the reference image
tagged \texttt{m6-memory-bus:thesis-v1.0}, with timestamps and chain
hashes elided for legibility.

\begin{verbatim}
# Step 1: Write a CONFIDENTIAL fragment with target ratio 4x.
$ curl -X POST http://localhost:8000/v1/write \
       -H "Content-Type: application/json" \
       -H "X-M6-Principal: principal=alice;acl=0xFFFF;cls=3" \
       -d @- <<'EOF'
{
  "fragment_id": "demo-001",
  "text": "Project P-0001 quarterly hours: 112; budget: EUR 18420.",
  "tags": {"acl_mask": 12876, "classification": 2},
  "target_ratio": 4.0,
  "compressor": "lingua2"
}
EOF
{"slot_id": "slot-1a4b9c2d", "audit_row": 1,
 "chain_hash": "5f3d2a1b...",  "achieved_ratio": 3.72}

# Step 2: Read the slot back. The principal's classification (3)
#         is at least the slot's (2) and the principal's ACL mask
#         is a superset, so the read is permitted.
$ curl -H "X-M6-Principal: principal=alice;acl=0xFFFF;cls=3" \
       http://localhost:8000/v1/read/slot-1a4b9c2d
{"slot_id": "slot-1a4b9c2d",
 "compressed_text": "Project P-0001 hours: 112 budget: 18420.",
 "compressor": "lingua2", "achieved_ratio": 3.72,
 "tags": {"acl_mask": 12876, "classification": 2},
 "created_at": "2026-05-30T07:42:15Z"}

# Step 3: A second principal with insufficient classification is
#         denied; the denial is itself audited with a
#         per-(principal, slot) 60-second deduplication.
$ curl -H "X-M6-Principal: principal=bob;acl=0xFFFF;cls=1" \
       -o /dev/null -w "%{http_code}\n" \
       http://localhost:8000/v1/read/slot-1a4b9c2d
403

# Step 4: Audit walk for the slot. The chain is verifiable end-to-
#         end via the SHA-256 chain hash on each row; tampering
#         with any row breaks verification.
$ curl -H "X-M6-Principal: principal=alice;acl=0xFFFF;cls=3" \
       http://localhost:8000/v1/audit/slot-1a4b9c2d
[
  {"rowid": 1, "event_type": "WRITE",  "result": "OK",
   "requester": "alice", "chain_hash": "5f3d2a1b..."},
  {"rowid": 2, "event_type": "READ",   "result": "OK",
   "requester": "alice", "chain_hash": "9c4e7a8d..."},
  {"rowid": 3, "event_type": "READ",   "result": "DENIED",
   "requester": "bob",   "chain_hash": "1b8a6f3e..."}
]
\end{verbatim}

The trace illustrates four properties: compression with achieved-
ratio reporting on the write path; policy-checked read with the
five-tier classification lattice and the access-control mask;
audited denial with the per-principal-per-slot deduplication that
prevents audit-log flooding; and chain-hash-verifiable provenance
on the audit endpoint. The complete OpenAPI schema is regenerated
by \texttt{make bus-openapi} and lives at \texttt{docs/openapi.json}
in the repository.

\subsection{D.3 Compute envelope and software stack}

The empirical work was carried out on the following stack.
The local workstation is an Apple Mac with the M4~Pro SoC,
$48$~GB of unified memory, Python~3.12, MLX framework version
$0.21$, Ollama version $0.3$ with the Qwen-2.5, Phi-3-Mini,
and Llama-3.1 model weights pre-pulled, and the
\texttt{m6} package installed in editable mode from the
repository. The GPU host is a Windows desktop running WSL2
Ubuntu~22.04 with the NVIDIA RTX~5090 $32$~GB GPU, CUDA toolkit
\textsuperscript{$\ge$}~12.0, and the same Python and Ollama
versions. The Tailscale mesh network connects the local
workstation to the GPU host as host \texttt{gpu} on the
\texttt{100.70.160.59/32} private address. The memory-bus
reference service runs on either host via
\texttt{docker compose -f docker/docker-compose.yml up -d}; the image is Python~3.12-slim
with FastAPI, Uvicorn, SQLite, and the \texttt{m6} package.

\section{Appendix E. RAG Pipeline Placement (H3)}
\markboth{left}{\normalfont{Appendix E. RAG Pipeline Placement}}
\label{appendix:rag}
\label{sec:h3_cost_caveat}

H3 examined where compression should sit relative to retrieval. It
is reported here, rather than in the main text, because its
pre-registered prediction was not satisfied and the salvaged result
is a single-stack non-reproduction orthogonal to the coordination
cliff and the memory bus.

Three pipeline placements are compared
(Section~\ref{sec:rag_pipelines} of \charef{implementation}):
compress-first (P1, compress the corpus then index it),
retrieve-first (P2, index the full corpus then compress the
top-$k$), and a joint relevance-conditional router (P3). Each runs
under a storage-bounded regime (FAISS index $\le 100$~MB) and an
accuracy-bounded regime (retrieval recall@10 $\le 0.85$) on the C1
family-a generator (\texttt{results/h3\_final/}). The pre-registered
H3 predicate was a \emph{sign-flip} in the optimal pipeline between
the two regimes with a $\ge\!5$~pp F1 effect. It is \emph{not}
observed: P1 is preferred over P2 in \emph{both} regimes (by $3.2$
and $2.0$~pp F1, paired-bootstrap $95\%$ CIs $[1.84, 4.72]$ and
$[1.15, 2.85]$, both excluding zero), which does not reproduce the
post-retrieval-compression preference of the
LongLLMLingua~\cite{jiang2024longllmlingua} line of work on the
single \texttt{bge-large-en-v1.5} + FAISS-CPU stack tested.

\begin{table}[h]
\centering
\caption{H3 verdict table. Mean F1, per-query EUR cost, and
\emph{achieved} compression ratio ($\times$) by regime and pipeline
from \texttt{results/h3\_final/} and the cost-instrumented re-run
\texttt{results/h3\_eprice/} ($n=150$ workloads per cell). The
sign-flip predicate is not satisfied (P1 is above P2 in both
regimes); P3 leads both but only P1 actually compresses, so the
cross-pipeline F1 gap is not a matched-compression comparison.}
\label{tab:h3}
\small
\begin{tabular}{lrrr|rrr}
\hline
& \multicolumn{3}{c|}{Storage-bounded}
& \multicolumn{3}{c}{Accuracy-bounded} \\
Pipeline & F1 & EUR/q & $\times$ & F1 & EUR/q & $\times$ \\
\hline
P1 (compress$\to$retrieve) & $0.396$ & $2.72\times 10^{-5}$ & $7.49$ & $0.648$ & $3.29\times 10^{-5}$ & $2.11$ \\
P2 (retrieve$\to$compress) & $0.364$ & $2.82\times 10^{-5}$ & $1.00$ & $0.628$ & $3.00\times 10^{-5}$ & $1.00$ \\
P3 (joint, conditional)    & $\mathbf{0.820}$ & $3.07\times 10^{-5}$ & $1.00$ & $\mathbf{0.895}$ & $3.10\times 10^{-5}$ & $1.00$ \\
\hline
$p_\text{Holm}$ (P1 vs P2, paired bootstrap) & \multicolumn{3}{c|}{$2\!\times\!10^{-4}$} & \multicolumn{3}{c}{$2\!\times\!10^{-4}$} \\
\hline
\end{tabular}
\end{table}

The joint router P3 leads both regimes on F1
(Table~\ref{tab:h3}), but a cost-instrumented re-run
(\texttt{results/h3\_eprice/}) shows it achieves this by routing
retrieved fragments \emph{verbatim} (achieved compression
$1.00\times$) while only P1 actually compresses ($7.49\times$
storage, $2.11\times$ accuracy). The \texttt{eur\_per\_query} cost
model is corpus-dominated (the three pipelines differ by only
$4.1$--$4.9\%$ in EUR within a regime) and, more importantly, does
\emph{not} meter compression compute at all, only the resulting
token counts, so a pipeline that runs an expensive compressor pass
is not charged for it. The thesis therefore makes \textbf{no
cost-parity and no Pareto claim} for the P3-vs-P1 comparison: the
cross-pipeline F1 gap is not a matched-compression comparison, and
P3 should be read as ``verbatim routing preserves F1 at an unmetered
compression cost''. A faithful per-pipeline cost model that meters
the compressor pass and amortises a persistent index across queries,
together with a multi-retriever replication, is the natural
follow-on (Section~\ref{sec:future}).
